% ---- pg-028 （书页 16） ----

\begin{mvfigure}
\mvinc{pg-028}{0.28277,0.12806,0.71298,0.73289}
\phantomsection\label{fig:1.10}\mvcap{图 1.10\quad 纸张性能数据的改进散布图}
{\zihao{6} (a) 选择离群值(25); (b) 删除离群值(25)}
\end{mvfigure}

行直到获得足够的信息.关于动态图解法的全面讨论可参见文献[1].关于与主机联在一起工作的多元探测图解分析的一个策略参见[32],该策略的动机是为了得到一个可寻找多元数据结构的常规程序.

\mvbold{例 1.9} (在三维空间作图)

第 4 章中的\hyperref[tab:4.3]{表 4.3} 给出了木材硬度的四个不同测量值.在例 4.14 中,样本(木板)16 及可能样本(木板)9 被标为异常观测值.\hyperref[fig:1.12]{图 1.12}(a),(b)和(c)显示出硬度数据在 $x_1,x_2,x_3$ 空间上%
